%%%%%%%%%%%%%%%%%%%%%%%%%%%%%%%%%%%%%%%%%%%%%%%%%%%%%%%%%%%%%%%%%%%%%%%%%%%%%%%%
% Official RA-L final-version style.
% RA-L final accepted manuscripts use IEEE journal style:
% \documentclass[letterpaper, 10 pt, journal, twoside]{IEEEtran}
%%%%%%%%%%%%%%%%%%%%%%%%%%%%%%%%%%%%%%%%%%%%%%%%%%%%%%%%%%%%%%%%%%%%%%%%%%%%%%%%

\documentclass[letterpaper, 10 pt, journal, twoside]{IEEEtran}

\usepackage{amsmath,amssymb}
\usepackage{graphicx}
\usepackage{booktabs}
\usepackage{xcolor}

\newcommand{\todo}[1]{\textcolor{red}{TODO: #1}}
\newcommand{\method}{ToolFrame}
\newcommand{\placeholder}[1]{\textit{#1}}

% Paper headers. Use only for final RA-L version.
\markboth{IEEE Robotics and Automation Letters. Preprint Version. Accepted
\placeholder{Month, Year}}
{FirstAuthorSurname \MakeLowercase{\textit{et al.}}: ToolFrame}

\title{\method{}: Let the Tool Define the Frame for Dexterous Manipulation}

\author{First Author$^{1}$, Second Author$^{2}$, and Third Author$^{1}$%
\thanks{Manuscript received: \placeholder{Month Day, Year}; Revised:
\placeholder{Month Day, Year}; Accepted: \placeholder{Month Day, Year}.}%
\thanks{This paper was recommended for publication by Editor
\placeholder{Editor A. Name} upon evaluation of the Associate Editor and
Reviewers' comments. This work was supported by
\placeholder{organizations/grants}.}%
\thanks{$^{1}$First Author and Third Author are with \placeholder{Department,
University, City, Country}. {\tt\footnotesize \placeholder{email}}}%
\thanks{$^{2}$Second Author is with \placeholder{Department, University, City,
Country}. {\tt\footnotesize \placeholder{email}}}%
\thanks{Digital Object Identifier (DOI): see top of this page.}%
}

\begin{document}
\maketitle

\begin{abstract}
The central difficulty in dexterous tool use is that success is determined by
an unobserved functional mechanism rather than by object identity alone: a
brush sweeps through bristle contact, an eraser wipes through a flat patch, a
hammer acts through an impact face, and a screwdriver acts through a tip and
shaft axis. A robot therefore must recover both the tool-side mechanism
(\emph{where} the tool should contact and \emph{which axis} should drive the
effect) and the scene-side mechanism (\emph{where} the effect should start,
move, and terminate). Directly asking a VLM to output robot trajectories hides
this mechanism inside free-form actions and makes evaluation vulnerable to
oracle axes, hidden target references, or task-family shortcuts. This paper
proposes \method{}, a tool-centric image-grounded pipeline that exposes the
mechanism as explicit intermediate references. The VLM is restricted to
selecting sparse public evidence: tool-side functional references
\texttt{tool.contact\_point} and \texttt{tool.motion\_axis}, and scene-side
RGB-D references for start, path, end, and direction. A geometry-constrained
compiler then binds the tool mechanism to the scene mechanism and solves
whole-tool $SE(3)$ subgoals, which are executed by a SimToolReal fixed-goal
dexterous policy. Task contracts evaluate the final physical effect and
attribute failures to functional grounding, scene grounding, subgoal validity,
policy execution, or task effect. \todo{Insert final headline results from
frozen launcher\_summary files and corresponding per-task summary.json files.}
\end{abstract}

\begin{IEEEkeywords}
Dexterous manipulation, robotic tool use, visual grounding, tool affordance,
motion planning.
\end{IEEEkeywords}

\section{Introduction}
\IEEEPARstart{A}{ tool} is useful because its geometry induces a task
mechanism. The robot does not merely need to know that an object is a brush,
eraser, hammer, screwdriver, or marker; it must know which part of the object
transfers force or contact to the scene and which motion direction produces the
intended effect. This is especially important for dexterous manipulation, where
the hand may hold the same object in many feasible ways, but only a small
subset of whole-tool poses activates the correct functional part.

The challenge is to recover this mechanism from public evidence without
collapsing it into an opaque policy or prompt-generated motion. If a VLM is
asked to generate dense waypoints or robot trajectories, the paper can no
longer cleanly answer whether success came from visual grounding, geometric
reasoning, controller robustness, or hidden task priors.

\method{} decomposes the mechanism into two sparse reference sets. The
tool-side set describes how the tool can create an effect: the functional part,
contact point, and motion axis. The scene-side set describes where the effect
should occur: start, path, end, and direction references lifted from RGB-D. A
compiler then solves the missing bridge between them: whole-tool $SE(3)$ poses
that make the tool mechanism coincide with the scene mechanism while respecting
clearance and stage constraints. In short, \emph{the tool defines the frame;
the hand follows}.

The paper contributes: (i) a mechanism-explicit, public-evidence tool-use
pipeline for dexterous manipulation; (ii) a tool-centric reference interface
based on \texttt{tool.contact\_point} and \texttt{tool.motion\_axis}; (iii) a
geometry-constrained compiler from tool and scene refs to whole-tool $SE(3)$
subgoals; and (iv) a task-contract benchmark protocol for final success and
failure attribution.

\section{Related Work}
\subsection{Vision-Language Models for Robot Planning}
Discuss VLM/LLM robot planning, visual grounding, language-conditioned
manipulation, and code-as-policies. The underlying challenge is that language
models can encode semantic priors but do not by themselves guarantee metric
contact geometry. The contrast is that \method{} uses VLMs to expose sparse
evidence for the mechanism, not to generate the mechanism's robot trajectory.

\subsection{Tool Affordances and Functional Part Grounding}
Discuss affordance learning, part segmentation, functional region prediction,
and contact-point estimation. The underlying challenge is that a functional
part label is not yet an executable manipulation constraint. The contrast is
that \method{} turns functional parts into explicit contact-point and
motion-axis constraints.

\subsection{Constraint-Based Manipulation and Tool-Use Benchmarks}
Discuss task-space constraints, contact-rich manipulation, keyframe/subgoal
planning, $SE(3)$ optimization, and tool-use benchmarks. \method{} binds tool
mechanism and scene mechanism before policy rollout, and task contracts
separate final task effect from first-failure attribution.

\section{Problem Formulation and Evidence Boundary}
\subsection{Mechanistic View of Tool Use}
We formulate a tool-use task as a relation between a tool mechanism and a scene
mechanism. The tool mechanism specifies a contact point and a motion axis in
the tool frame. The scene mechanism specifies the spatial locus and direction
of the desired effect in the world frame. The robot succeeds only if the
executed whole-tool motion makes these two mechanisms coincide under the task
contract.

\subsection{Allowed and Forbidden Inputs}
The method may use natural task language, public RGB-D scene evidence, camera
calibration, public tool meshes, segmented tool-part evidence, reviewed or
review-pending public tool references, and robot state needed for execution. It
must not use hidden evaluator target references, simulator target centers,
oracle target axes, oracle trajectories, tool-family fallback, task-family
fallback, primitive fallback, repair-loop fallback, or seed-specific shortcuts.

\subsection{Reference Interface}
Tool references are
\begin{equation}
  \mathcal{R}_{tool} =
  \{p^{contact}_{tool}, a^{motion}_{tool}\},
\end{equation}
and scene references are
\begin{equation}
  \mathcal{R}_{scene} =
  \{p^{start}_{scene}, p^{path}_{scene}, p^{end}_{scene},
  a^{direction}_{scene}\}.
\end{equation}
The compiler solves whole-tool poses $T_{world,tool}\in SE(3)$ that align the
tool contact and motion references with the grounded scene references.

\section{Method}
\subsection{Overview}
\begin{figure*}[t]
\centering
\fbox{\parbox{0.92\textwidth}{\centering
Placeholder for Fig. 1: (A) Tool-Centric Prior Extraction;
(B) Scene Key-Points Grounding; (C) Geometry-Constrained Optimizer;
(D) Policy Rollout and Evaluation.}}
\caption{\method{} overview. Public tool-part evidence exposes the tool
mechanism; public RGB-D evidence exposes the scene mechanism; a tool-centric
$SE(3)$ compiler binds both into whole-tool subgoals executed by a SimToolReal
fixed-goal dexterous policy.}
\label{fig:overview}
\end{figure*}

\subsection{Tool-Centric Prior Extraction}
The challenge on the tool side is that visual appearance and object category do
not identify the operational mechanism. This module extracts the functional
part, \texttt{tool.contact\_point}, and \texttt{tool.motion\_axis} from public
tool evidence. The output is the tool mechanism that later constrains
whole-tool poses.

\subsection{Scene Key-Points Grounding}
The challenge on the scene side is that the task effect is spatially sparse.
\method{} uses the VLM only to select public RGB-D start, path, end, and
direction refs. These 2D references are lifted to 3D using depth and camera
calibration.

\subsection{Geometry-Constrained Optimizer}
The compiler solves for whole-tool $SE(3)$ poses that satisfy contact-point
matching, motion-axis alignment, clearance, pre-contact approach, and staged
progression. The output is a sequence of fixed whole-tool subgoals, not
low-level actions or VLM-generated trajectories.

\subsection{Policy Rollout and Evaluation}
The SimToolReal fixed-goal policy executes the compiled whole-tool subgoals.
Task contracts evaluate the resulting physical effect after rollout and
attribute the first failure to the responsible stage. Contract outcomes are not
used as a repair loop or fallback during the method.

\section{Benchmark and Experiments}
\begin{table}[t]
\caption{Functional tool geometries targeted by \method{}.}
\label{tab:tools_final}
\centering
\begin{tabular}{ll}
\toprule
Tool & Functional mechanism \\
\midrule
Brush & broad contact and sweep axis \\
Eraser & flat patch and wipe axis \\
Hammer & impact face and strike axis \\
Screwdriver & tip and insertion axis \\
Marker & point tip and writing axis \\
\bottomrule
\end{tabular}
\end{table}

The benchmark should stress whether the mechanism is genuinely geometric rather
than scene-specific, covering tabletop and side-board scene families. Headline
statistics must be reported only for tasks with current/e2e contracts and
frozen summaries. Candidate, diagnostic, and draft/smoke tasks should be
labeled separately.

The main results table should report functional grounding rate, scene grounding
rate, subgoal validity rate, policy execution rate, final task success rate,
and failure attribution. \todo{Fill only from frozen launcher summaries and
matched per-task summaries.}

\section{Limitations and Conclusion}
Mechanism recovery requires visible RGB-D evidence and accurate 2D-to-3D
lifting. Tool-side mechanisms depend on public tool-part evidence and human
review. Valid mechanism alignment can still fail under the fixed-goal policy.
The system intentionally does not use online repair, primitive fallback, or
seed-specific shortcuts.

\method{} treats dexterous tool use as the recovery and execution of a
functional mechanism. The tool supplies the contact and motion frame, the scene
supplies the desired effect frame, and the compiler binds the two into
whole-tool $SE(3)$ subgoals for fixed-goal dexterous execution.

\bibliographystyle{IEEEtran}
\bibliography{references}

\end{document}
